\documentclass[journal]{IEEEtran}

\ifCLASSINFOpdf
\else
   \usepackage[dvips]{graphicx}
\fi
\usepackage{url}
\hyphenation{op-tical net-works semi-conduc-tor}
\usepackage{graphicx}
\usepackage{amsmath}
\usepackage{algorithm} 
\usepackage{algorithmic} 
\usepackage{booktabs}
\usepackage{cite}
\usepackage{xcolor}
\usepackage{multirow}
\usepackage[normalem]{ulem}
\useunder{\uline}{\ul}{}
\usepackage{graphicx}
\usepackage{subcaption}
\usepackage{arydshln}


\begin{document}

\title{Connectivity-Aware Client Selection in Decentralized Federated Learning}

\author{Usevalad Milasheuski, \IEEEmembership{Member, IEEE}, Monica Nicoli, \IEEEmembership{Senior Member, IEEE}, and Stefano Savazzi, \IEEEmembership{Senior Member, IEEE}
\thanks{This paragraph of the first footnote will contain the date on which you submitted your paper for review. It will also contain support information, including sponsor and financial support acknowledgment. For example, ``This work was supported in part by the U.S. Department of Commerce under Grant BS123456.'' }
\thanks{The next few paragraphs should contain the authors' current affiliations, including current address and e-mail. For example, F. A. Author is with the National Institute of Standards and Technology, Boulder, CO 80305 USA (e-mail: author@boulder.nist.gov).}
\thanks{S. B. Author, Jr., was with Rice University, Houston, TX 77005 USA. He is now with the Department of Physics, Colorado State University, Fort Collins, CO 80523 USA (e-mail: author@lamar.colostate.edu).}}

\markboth{Journal of \LaTeX\ Class Files, Vol. 14, No. 8, August 2015}
{Shell \MakeLowercase{\textit{et al.}}: Bare Demo of IEEEtran.cls for IEEE Journals}
\maketitle

\begin{abstract}
Decentralized federated learning (DFL) has emerged as a scalable alternative to centralized federated learning, removing the single point of failure and communication bottleneck at the server. However, most existing decentralized approaches rely on full-neighborhood broadcasting, where each client transmits its model update to all of its neighbors in each communication round. While effective, this results in significant communication overhead and large energy consumption. This paper proposes novel client selection methods that leverage the topology of the nodes by analyzing the spectral properties of the graph Laplacian. Specifically, we exploit eigenvector embeddings and heat/diffusion kernel similarities to guide each client in selecting only a subset of structurally relevant neighbors, i.e. those that are most beneficial for the convergence of the federation process. This topology-aware randomization is shown to outperform uniform neighborhood subsampling, while reducing the total number of transmitted messages compared to naive broadcasting.
\end{abstract}

\begin{IEEEkeywords}
Decentralized Federated Learning, Graph Spectral Methods, Client Selection, Heat/Diffusion Kernel
\end{IEEEkeywords}


\IEEEpeerreviewmaketitle



\section{Introduction}

\IEEEPARstart{F}{ederated} Learning (FL) is a distributed paradigm that enables collaborative training of a global model across decentralized edge devices without sharing raw data \cite{FL}. Unlike centralized training, where data are aggregated on a single server, FL performs on-device training and transmits only model updates. This preserves privacy, reduces communication costs, and makes FL a promising solution for diverse applications. 

While FL provides significant privacy benefits, its reliance on a central server for model aggregation introduces a single point of failure. The entire training process is susceptible to server-side issues, including hardware failures, network outages, or malicious attacks, which can halt the process and compromise the global model's integrity \cite{attack1, attack2}. To overcome this vulnerability, a paradigm known as Decentralized Federated Learning (DFL) has emerged. DFL eliminates the central server, instead leveraging a peer-to-peer network where devices communicate directly with their neighbors to share and aggregate model updates. This distributed architecture not only removes the single point of failure but also improves scalability and resilience, as the system can continue to operate even if some devices fail or become disconnected \cite{decSGDKOLOS, DFLsurvey, DecSGDTop}.


In many studies of DFL, a Complete connected communication graph is assumed, where each client transmits its model updates to all other clients \cite{DSGD, healthcare, decOPT, iot}. This all-to-all communication model simplifies convergence analysis and provides strong theoretical guarantees, as it ensures that every client's model is quickly and thoroughly updated by the entire network. This approach guarantees information propagation but results in high communication overhead, which is particularly problematic in bandwidth-constrained or energy-limited environments \cite{decOPT}. To overcome this limitation, more practical DFL approaches focus on designing sparse communication topologies or employing intelligent client selection strategies \cite{graph_select1}. These methods aim to reduce the communication burden while still maintaining satisfactory convergence rates and model performance. Most existing works construct the communication graph in one of three ways: (i) from an initially Complete graph using a similarity measure \cite{smart_samp, select_data_2}, (ii) by selecting neighbors based on model or gradient similarity \cite{select_grad_1}, or (iii) by considering divergences in the quality or distribution of local heterogeneous datasets \cite{select_data_2}. However, these methods ignore the fixed communication topology, which restricts connections between certain pairs of nodes.

In this paper, we address this inefficiency by \textit{connectivity-aware neighbor selection}. Instead of naive broadcasting, clients select neighbors according to similarity measures derived from the graph Laplacian of the adjacency matrix of the clients. Our goal is to propose decentralized client selection methods, which rely only on the topology of the clients and are more efficient than broadcasting in terms of energy and communication, without analyzing data distribution or model parameter dynamics. In particular, leverage two approaches:

\begin{itemize} 
    \item \textbf{Spectral Embedding}: representing clients in a low-dimensional space using Laplacian eigenvectors and selecting neighbors that are close in this embedding. 
    \item \textbf{Heat Kernel Similarity}: interpreting diffusion over the graph as a similarity measure and selecting neighbors with the highest heat kernel weights. 
\end{itemize}
The paper is organized as follows. Section \ref{sec:relwork} reviews the related work. Section \ref{sec:methodology} discusses the proposed topology-aware client selection methods. Results are validated in Section \ref{sec:exper} with real world datasets. Concluding remarks are in Section \ref{sec:concl}.
 


\section{Related Work and Contributions}
\label{sec:relwork}

DFL enables collaborative training without a central server, protecting data privacy, and is used in drones \cite{drones}, IoT \cite{iot}, and healthcare \cite{healthcare}. However, decentralized methods often assume Complete topologies and broadcasting \cite{DSGD}, which can make communication inefficient despite strong convergence guarantees.

For example, the naive centralized FedAvg algorithm \cite{FL} randomly selects a subset of clients each round to address this issue. Some methods use importance sampling \cite{imp_samp1, imp_samp2}, prioritizing clients with large gradient norms, while others measure the gradient or model alignment between models \cite{client_GRAD_sampling, client_GRAD_sampling2} to effectively subsample the neighborhood. In decentralized setups, a new graph is often constructed over the Complete graph topology by specifying the adjacency matrix or learning implicit similarities from data variability, update norms, or validation accuracies \cite{graph_select1, graph_select2}. While effective under heterogeneous data, decentralized methods ignore predefined client topologies, causing bottlenecks and inefficient information flow. Graph-theoretic techniques represent nodes and assess their importance via metrics like degree, centrality, clustering coefficient, resource allocation index, or random-walk methods such as DeepWalk \cite{deepwalk} and Node2Vec \cite{node2vec}.

A particularly powerful and richer class of methods for graph representation involves spectral graph theory \cite{graph_theory}. Unlike simple ad-hoc methods, spectral methods analyze the graph's properties through the eigenvalues and eigenvectors of its adjacency or Laplacian matrix. The eigenvectors of the graph Laplacian encode a unique representation of each node's position and its influence on the network's overall connectivity and information flow. This can allow for a more nuanced approach to client selection, where nodes can be prioritized based on their spectral properties to optimize the global aggregation process. In this way, we can ensure a more rapid and stable model convergence, while dramatically reducing the communication burden, thus making DFL a more scalable and reliable solution. These ideas are used for multiple applications, including graph partitioning, node clustering \cite{graph_theory} and positional encoding for transformers \cite{pos_emb}.

\section{Methodology}
\label{sec:methodology}

This section introduces the proposed methods for neighbor selection based on the spectral properties of the topology graph of decentralized clients. The goal is to efficiently select the most effective nodes to federate with out of the entire neighborhood to reduce the communication overhead and improve model convergence and generalization.

\subsection{Problem Setup} 
We consider an undirected  graph $G=(V,E)$ where $\mathcal{V} = \{1,2,\dots,N\}$ is the set of $N$ nodes/clients and $\mathcal{E} \subseteq \mathcal{V} \times \mathcal{V}$ is the set of communication links. The adjacency matrix $\mathbf{A} \in \{0,1\}^{N \times N}$ is defined as
\begin{equation}
    A_{ij} =
        \begin{cases}
            1, & \text{if } (i,j) \in \mathcal{E} \text{ and } i \neq j, \\
            0, & \text{otherwise.}
        \end{cases}
\end{equation}

The unnormalized Laplacian of a graph is defined as \cite{graph_theory}:
\begin{equation}
    \mathbf{L} = \mathbf{D} - \mathbf{A},
\end{equation}
where $\mathbf{D}$ is a diagonal degree matrix with entries $D_{ii} = \sum_{j=1}^N A_{ij}, \quad \forall i \in \{1, \dots, N\}$. For spectral analysis, a normalized version of the Laplacian is commonly used:

\begin{equation}
    \mathbf{L}_{sym} = \mathbf{D}^{-1/2}\mathbf{L}\mathbf{D}^{-1/2},
\end{equation}

Both Laplacians $\mathbf{L}$ and $\mathbf{L}_{sym}$ are symmetric and positive semi-definite. 
Hence, it admits the eigenvalue decomposition
\begin{equation}    
\mathbf{L}_{sym} = \mathbf{U} \mathbf{\Lambda} \mathbf{U}^\top,
\end{equation}
where $\mathbf{U} = [\mathbf{u}_1, \mathbf{u}_2, \dots, \mathbf{u}_N] \in \mathbf{R}^{N \times N}$ is an orthonormal matrix of eigenvectors, 
and $\mathbf{\Lambda} = \mathrm{diag}(\lambda_1, \lambda_2, \dots, \lambda_N)$ contains the eigenvalues 
ordered as
\begin{equation}
0 = \lambda_1 \leq \lambda_2 \leq \cdots \leq \lambda_N.
\end{equation}


The spectrum of the graph Laplacian reveals both structural and dynamical properties of the underlying communication network. The eigenvalues describe the "frequencies" of the graph: the smallest nonzero eigenvalue indicates how well the clients are linked, with larger values corresponding to faster consensus and more robust information flow \cite{graph_theory}. The eigenvectors can be understood as graph harmonics, where smooth eigenvectors (corresponding to small eigenvalues) vary slowly across connected clients, while those tied to larger eigenvalues fluctuate more rapidly, capturing local fluctuations in the topology. These properties will be exploited in the following sections for the subsampling of the neighborhood.

\subsection{Spectral Embedding} 
Laplacian eigenvectors form the basis of spectral embeddings that map each client into a low-dimensional Euclidean space. Let $\mathbf{U}_k \in \mathbf{R}^{N \times k}$ contain the $k$ eigenvectors $\mathbf{u}_i=[u_{i,n}]_{n=1}^{N}$, $i=2, ...,k+1$ corresponding to $k$ smallest nonzero eigenvalues, the embedding $\mathbf{e}_n$ of node $n$ is given by the normalized $n$-th \textit{row} of $\mathbf{U}_k$:

\begin{equation} 
    \mathbf{e}_n = [\tilde{u}_{2,n}, \tilde{u}_{3,n}, \dots, \tilde{u}_{k+1,n}], 
\label{eq:embeddings}
\end{equation} 
where $\tilde{u}_{i,n} = u_{i,n}/(\sum_{i=2}^{k+1} u_{i,n}^2 )^{1/2}$. 
By selecting the $k$ smallest non-zero eigenvectors, the embedding from (\ref{eq:embeddings}) sets a sufficient representation of each individual node. This representation provides geometric coordinates for each client, where distances reflect connectivity patterns in the network, enabling deeper analysis of decentralized learning dynamics. 

Taking inspiration from node embedding techniques \cite{deepwalk, node2vec}, we can compare the similarity between the nodes based on the embedding from (\ref{eq:embeddings}) as:
\begin{equation} \label{eq:cos_sim}
    S^{spectr}(n_1, n_2) = \sum_{i=2}^{k+1}  \tilde{u}_{i,n_1} \cdot \tilde{u}_{i,n_2}.
\end{equation}
The cosine similarity between two embedded vectors provides a measure of how aligned their positions are with respect to the network's global connectivity structure. A high cosine similarity indicates that the two clients occupy similar “directions”, suggesting they share comparable communication neighborhoods and thereby play analogous roles in information diffusion. Conversely, a low similarity implies that the clients lie in orthogonal directions in the spectral space, reflecting structural dissimilarity in the graph.

\subsection{Heat Kernel}
The spectral properties of the Laplacian also allow the construction of diffusion-based similarity measures. A central example is the heat kernel, which models the propagation of information across the communication graph in analogy with heat diffusion over a surface \cite{heat_kern}. Formally, the heat kernel at diffusion time $t > 0$ is defined as

\begin{equation} \label{eq:dif_ker}
    \mathbf{K}_t = \exp(-t\mathbf{L}) = \sum_{n=1}^N e^{-t\lambda_n}\mathbf{u}_n \mathbf{u}_n^T, 
\end{equation}
where $t$ is temperature: small values of $t$ emphasize local structure, highlighting immediate neighborhoods, while larger values of $t$ reflect long-range connectivity and global mixing behavior. Each entry $K_{n_1 n_2}(t)$ captures the amount of “heat” transferred from node $n_1$ to node $n_2$ after time $t$. We can then use inverse diffusion distance \cite{heat_kern} to quantify how information diffuses between the two clients, and use it as a similarity between the nodes, as an alternative to Eq. \ref{eq:cos_sim}:
\begin{equation}\label{eq:diff_dist}
    S^{heat}(n_1, n_2) =  (K_{n_1 n_1}(t) + K_{n_2 n_2}(t) - 2K_{n_1 n_2}(t))^{-1}.
\end{equation}


\subsection{Decentralized Federated Learning}

We consider the set of clients over the previously defined graph $G$. Each client $n$ holds local data $\mathcal{D}^n=\{\mathbf{x}^n_i, y^n_i\}_{i=1}^{|\mathcal{D}^n|}$, where $\mathbf{x}^n_i$ and $y^n_i$ represent local features and labels of client $n$ respectively. Generally, the data between the clients does not follow the Independent and Identically Distributed (IID) assumption, leading to highly heterogeneous setups. The goal is to collaboratively find a set of model parameters $\boldsymbol{w}$ that minimize the average loss over the nodes

\begin{equation} \label{eq:optim}
    \min_{\boldsymbol{w} \in \mathbf{R}^d} L(\boldsymbol{w})=\min_{\boldsymbol{w} \in \mathbf{R}^d} \frac{1}{N}\sum_{n=1}^{N}L_n(\boldsymbol{w}),
\end{equation}
where  $L_n(\boldsymbol{w})=\frac{1}{|\mathcal{D}^n|}\sum_{(\mathbf{x}^n_i, y^n_i) \in \mathcal{D}^n}l(\boldsymbol{w}, (\mathbf{x}^n_i, y^n_i))$. 

Given a node $n$ and its neighbors $\mathcal{N}(n)$, the objective (\ref{eq:optim}) is minimized iteratively. At each round $t$ the nodes: \textbf{1)} train model on the local dataset; \textbf{2)} forward the weights to the neighbors and \textbf{3)} aggregate the received models: 

\begin{equation} \label{eq:aggreg}
    \boldsymbol{w}_n^{t+1} = \rho  \boldsymbol{w}_n^{t} + (1 - \rho) \frac{1}{|\mathcal{N}(n)|}\sum_{m \in \mathcal{N}(n)} \boldsymbol{w}_m^{t},
\end{equation}
where $\rho$ controls the stability of the update.

%The key modification in our framework is that for each node $n$ instead of sending updates to all $|\mathcal{N}(n)|$ neighbors, as in (\ref{eq:aggreg}), we propose that the client samples from the neighborhood $\lceil \beta|\mathcal{N}(n)| \rceil$ neighbors, 
Instead of sending updates to all $|\mathcal{N}(n)|$ neighbors, as in (\ref{eq:aggreg}), in the proposed framework the client samples from a neighborhood $\lceil \beta|\mathcal{N}(n)| \rceil$, where $0\leq\beta\leq 1$ is the fraction of neighbors sampled at each node, according to a Laplacian-based similarity from (\ref{eq:cos_sim}) or (\ref{eq:diff_dist}):

\begin{equation} \label{eq:proba}
    p(n_1, n_2) = \frac{\exp({S(n_1, n_2)})}{\sum_{n \in \mathcal{N}(n_1)} \exp(S(n_1, n))}.
\end{equation}
For the graph in Figure \ref{fig:res}, we show an example of how these similarity metrics affect the transmission probability between a pair of connected nodes. As we can see, both approaches adjust the selection probabilities in a consistent way, which is evident from the varying colors of a node’s outgoing edges, indicating a non-uniform distribution.


\begin{figure*}[t]
    \centering
    % First subfigure
    \begin{subfigure}[b]{0.42\textwidth}
        \centering
        \includegraphics[width=\linewidth]{figures/small_spectr.png}
        \caption{Spectral Embedding}
        \label{fig:res3}
    \end{subfigure}
    \hfill
    % Second subfigure
    \begin{subfigure}[b]{0.42\textwidth}
        \centering
        \includegraphics[width=\linewidth]{figures/small_heatkern.png}
        \caption{Heat Kernel}
        \label{fig:res4}
    \end{subfigure}
    
    \caption{Selection probabilities for \textit{Spectral Clustering} (left) and \textit{Heat Kernel} methods on a \textit{small} graph.}
    \label{fig:res}
\end{figure*}


\section{Experiments}
\label{sec:exper}

We evaluate the performance of the proposed methods on a small-scale synthetic Erdős-Rényi graph \cite{random_graph} with 8 nodes and an edge probability of $0.4$, which we will address as \textbf{small} set, and a larger real-world network, the Davis Southern women social network graph \cite{south_graph} with $32$ nodes, which we will call \textbf{large} set. Two datasets have been used: FMNIST and CIFAR10 under the heterogeneous setup simulated using symmetrical Dirichlet distribution with $\alpha=0.3$ \cite{dir}.

We employ the models described in \cite{FL}, with the sole modification that FMNIST is used in place of MNIST, while CIFAR10 remains unchanged. The models use standard SGD for $3$ local epochs with a learning rate of $0.01$, a batch size of $64$ and $\rho=0.9$. We run the experiments for $100$ rounds and average the performance over $3$ runs. For the client selection, we use $\beta=0.5$, $k=3$ for the small graph and $k=12$ for the large graph for spectral embedding, $t=0.01$ and $t=0.05$ for the respective graphs for the heat kernel.  


\subsection{Experimental Setups}
The current work analyzes the way we can efficiently select a subset of each neighborhood by relying only on the connectivity graph, ignoring local datasets and model parameters. Thus, we compare the following setups, including baselines: 1) \textbf{No-federation} - average performance of independent local optimizations on individual datasets; 2) \textbf{Broadcasting} - each client broadcasts its model to all the nodes in the neighborhood; 3) \textbf{Random Sampling} - a random selection of $\beta \%$ of the neighboring nodes; 4) \textbf{Gradient Similarity} - selection of clients is performed by computing the angle between the normalized gradients of the node and the neighbors \cite{client_GRAD_sampling}; 5) \textbf{Spectral Embedding} - selection of $\beta \%$ of the neighbors based on (\ref{eq:cos_sim}) and (\ref{eq:proba}); 6) \textbf{Heat Kenel} - selection of $\beta \%$ of the neighbors based on (\ref{eq:dif_ker}) and (\ref{eq:proba}).




\subsection{Results} 
\label{sec:res}
The results in Table \ref{tab:res_table} highlight clear patterns across datasets and graph sizes. The \textit{no-federation} baseline yields the lowest performance, while \textit{broadcasting} achieves the best accuracy at the cost of maximal communication. Our topology-aware subsampling methods provide a middle ground: they consistently outperform random selection while remaining more communication and energy efficient than broadcasting.

On FMNIST, \textit{spectral embedding} emerges as the strongest subsampling method across both small and large graphs, showing that structural information can be effectively leveraged in simpler tasks. On the more challenging CIFAR10 dataset, performance differences are smaller but still meaningful: the \textit{heat kernel} method performs best on small graphs, while \textit{spectral embedding} is slightly better on large graphs. These results suggest that the relative benefits of each method depend on both dataset complexity and network size.

Finally, we observe that larger graphs generally reduce the overall accuracy due to increased heterogeneity, but our methods maintain a consistent advantage over random sampling. This robustness indicates that topology-aware client selection narrows the gap to the broadcast upper bound while keeping communication costs manageable, making it a practical solution for scalable decentralized learning.


\begin{table}[]
\centering
\begin{tabular}{|c|c|clcl|}
\hline
\multirow{2}{*}{Dataset} & \multirow{2}{*}{Selection} & \multicolumn{4}{c|}{Topology}                                            \\ \cline{3-6} 
                         &                            & \multicolumn{2}{c|}{Small}         & \multicolumn{2}{c|}{Large}          \\ \hline
\multirow{5}{*}{FMNIST}  & \textcolor{gray}{Broadcast}                  & \multicolumn{2}{c}{\textcolor{gray}{86.36}} & \multicolumn{2}{c|}{\textcolor{gray}{84.13}} \\ 
                        
                         & No-Fed.                    & \multicolumn{2}{c}{65.65}          & \multicolumn{2}{c|}{61.49}          \\
                         & Random                     & \multicolumn{2}{c}{81.01}          & \multicolumn{2}{c|}{80.84}          \\
                         & Grad. Sim.                    & \multicolumn{2}{c}{\textbf{82.53}}          & \multicolumn{2}{c|}{81.36}          \\
    
                         & Spectr. Emb.               & \multicolumn{2}{c}{ \underline{82.44}}    & \multicolumn{2}{c|}{\textbf{82.15}}    \\
                         & Heat Kern.                 & \multicolumn{2}{c}{81.31}          & \multicolumn{2}{c|}{\underline{81.86}}          \\ \hline
\multirow{5}{*}{CIFAR10} & \textcolor{gray}{Broadcast}                  & \multicolumn{2}{c}{\textcolor{gray}{62.02}} & \multicolumn{2}{c|}{\textcolor{gray}{54.53}} \\
        
                         & No Fed.                    & \multicolumn{2}{c}{40.68}          & \multicolumn{2}{c|}{31.53}          \\
                         & Random                     & \multicolumn{2}{c}{57.31}          & \multicolumn{2}{c|}{48.87}          \\
                         & Grad. Sim.                    & \multicolumn{2}{c}{57.5}          & \multicolumn{2}{c|}{48.64}          \\
    
                          & Spectr. Emb.               & \multicolumn{2}{c}{\underline{57.74}}          & \multicolumn{2}{c|}{\textbf{50.13}}  \\
                         & Heat Kern.                 & \multicolumn{2}{c}{\textbf{58.05}}    & \multicolumn{2}{c|}{\underline{49.89}}          \\ \hline
\end{tabular}
\caption{Average model accuracy across datasets and client topologies. The \textbf{bold values} indicate the best subsampling method  (excluding the
Broadcast), while the \underline{underlined values} highlight the second best approach.} % Best results are in bold and the second best are underscored.
\label{tab:res_table}
\end{table}

\section{Conclusion} 
\label{sec:concl}
We introduced a topology-aware client selection strategy for decentralized federated learning. By leveraging Laplacian spectral embeddings and diffusion kernels, clients can intelligently sample neighbors,  outperforming random selection and achieving faster convergence with fewer communications, leading to smaller transmission overhead and energy consumption, which can be critical for constrained applications. Future work will analyze how to combine multiple selection strategies in DFL based on the topology, data distribution and model gradients.

\bibliographystyle{IEEEtran}
\bibliography{bibliography}

\end{document}